% Original pages: 65-72
% Figures: none; one commutative diagram typeset directly
% Uncertainty log: none

\noindent\textbf{11.} If $\mathfrak p$ is a minimal prime ideal of a ring $A$, show that $S_\mathfrak p(0)$ (Exercise 10) is the
smallest $\mathfrak p$-primary ideal.

Let $\mathfrak a$ be the intersection of the ideals $S_\mathfrak p(0)$ as $\mathfrak p$ runs through the minimal
prime ideals of $A$.  Show that $\mathfrak a$ is contained in the nilradical of $A$.

Suppose that the zero ideal is decomposable.  Prove that $\mathfrak a=0$ if and only if
every prime ideal of 0 is isolated.

\medskip
\noindent\textbf{12.} Let $A$ be a ring, $S$ a multiplicatively closed subset of $A$.  For any ideal $\mathfrak a$, let $S(\mathfrak a)$
denote the contraction of $S^{-1}\mathfrak a$ in $A$.  The ideal $S(\mathfrak a)$ is called the \emph{saturation} of $\mathfrak a$
with respect to $S$.  Prove that

i) $S(\mathfrak a)\cap S(\mathfrak b)=S(\mathfrak a\cap\mathfrak b)$

ii) $S(r(\mathfrak a))=r(S(\mathfrak a))$

iii) $S(\mathfrak a)=(1)\Longleftrightarrow\mathfrak a$ meets $S$

iv) $S_1(S_2(\mathfrak a))=(S_1S_2)(\mathfrak a)$.

If $\mathfrak a$ has a primary decomposition, prove that the set of ideals $S(\mathfrak a)$ (where $S$ runs
through all multiplicatively closed subsets of $A$) is finite.

\medskip
\noindent\textbf{13.} Let $A$ be a ring and $\mathfrak p$ a prime ideal of $A$.  The \emph{$n$th symbolic power of $\mathfrak p$} is defined
to be the ideal (in the notation of Exercise 12)
\[
\mathfrak p^{(n)}=S_\mathfrak p(\mathfrak p^n)
\]
where $S_\mathfrak p=A-\mathfrak p$.  Show that

i) $\mathfrak p^{(n)}$ is a $\mathfrak p$-primary ideal;

ii) if $\mathfrak p^n$ has a primary decomposition, then $\mathfrak p^{(n)}$ is its $\mathfrak p$-primary component;

iii) if $\mathfrak p^{(m)}\mathfrak p^{(n)}$ has a primary decomposition, then $\mathfrak p^{(m+n)}$ is its $\mathfrak p$-primary component;

iv) $\mathfrak p^{(n)}=\mathfrak p^n\Longleftrightarrow\mathfrak p^{(n)}$ is $\mathfrak p$-primary.

\medskip
\noindent\textbf{14.} Let $\mathfrak a$ be a decomposable ideal in a ring $A$ and let $\mathfrak p$ be a maximal element of the
set of ideals $(\mathfrak a:x)$, where $x\in A$ and $x\notin\mathfrak a$.  Show that $\mathfrak p$ is a prime ideal belonging
to $\mathfrak a$.

\medskip
\noindent\textbf{15.} Let $\mathfrak a$ be a decomposable ideal in a ring $A$, let $\Sigma$ be an isolated set of prime ideals
belonging to $\mathfrak a$, and let $\mathfrak q_\Sigma$ be the intersection of the corresponding primary
components.  Let $f$ be an element of $A$ such that, for each prime ideal $\mathfrak p$ belonging
to $\mathfrak a$, we have $f\in\mathfrak p\Longleftrightarrow\mathfrak p\notin\Sigma$, and let $S_f$ be the set of all powers of $f$.  Show that
$\mathfrak q_\Sigma=S_f(\mathfrak a)=(\mathfrak a:f^n)$ for all large $n$.

\medskip
\noindent\textbf{16.} If $A$ is a ring in which every ideal has a primary decomposition, show that every
ring of fractions $S^{-1}A$ has the same property.

\medskip
\noindent\textbf{17.} Let $A$ be a ring with the following property.

(L1) For every ideal $\mathfrak a\ne(1)$ in $A$ and every prime ideal $\mathfrak p$, there exists $x\notin\mathfrak p$ such
that $S_\mathfrak p(\mathfrak a)=(\mathfrak a:x)$, where $S_\mathfrak p=A-\mathfrak p$.

Then every ideal in $A$ is an intersection of (possibly infinitely many) primary
ideals.

[Let $\mathfrak a$ be an ideal $\ne(1)$ in $A$, and let $\mathfrak p_1$ be a minimal element of the set of prime
ideals containing $\mathfrak a$.  Then $\mathfrak q_1=S_{\mathfrak p_1}(\mathfrak a)$ is $\mathfrak p_1$-primary (by Exercise 11), and $\mathfrak q_1=
(\mathfrak a:x)$ for some $x\notin\mathfrak p_1$.  Show that $\mathfrak a=\mathfrak q_1\cap(\mathfrak a+(x))$.

Now let $\mathfrak a_1$ be a maximal element of the set of ideals $\mathfrak b\supseteq\mathfrak a$ such that
$\mathfrak q_1\cap\mathfrak b=\mathfrak a$, and choose $\mathfrak a_1$ so that $x\in\mathfrak a_1$, and therefore $\mathfrak a_1\nsubseteq\mathfrak p_1$.  Repeat the
construction starting with $\mathfrak a_1$, and so on.  At the $n$th stage we have $\mathfrak a=\mathfrak q_1\cap\cdots
\cap\mathfrak q_n\cap\mathfrak a_n$ where the $\mathfrak q_i$ are primary ideals, $\mathfrak a_n$ is maximal among the ideals $\mathfrak b$ containing $\mathfrak a_{n-1}=\mathfrak a_n\cap\mathfrak q_n$ such that $\mathfrak a=\mathfrak q_1\cap\cdots\cap\mathfrak q_n\cap\mathfrak b$, and $\mathfrak a_n\nsubseteq\mathfrak p_n$.  If at any
stage we have $\mathfrak a_n=(1)$, the process stops, and $\mathfrak a$ is a finite intersection of
primary ideals.  If not, continue by transfinite induction, observing that each $\mathfrak a_n$
strictly contains $\mathfrak a_{n-1}$.]

\medskip
\noindent\textbf{18.} Consider the following condition on a ring $A$:

(L2) Given an ideal $\mathfrak a$ and a descending chain $S_1\supseteq S_2\supseteq\cdots\supseteq S_n\supseteq\cdots$
of multiplicatively closed subsets of $A$, there exists an integer $n$ such that $S_n(\mathfrak a)=
S_{n+1}(\mathfrak a)=\cdots$.  Prove that the following are equivalent:

i) Every ideal in $A$ has a primary decomposition;

ii) $A$ satisfies (L1) and (L2).

[For i) $\Rightarrow$ ii), use Exercises 12 and 15.  For ii) $\Rightarrow$ i) show, with the notation of
the proof of Exercise 17, that if $S_n=S_{\mathfrak p_1}\cap\cdots\cap S_{\mathfrak p_n}$ then $S_n$ meets $\mathfrak a_n$, hence
$S_n(\mathfrak a_n)=(1)$, and therefore $S_n(\mathfrak a)=\mathfrak q_1\cap\cdots\cap\mathfrak q_n$.  Now use (L2) to show that
the construction must terminate after a finite number of steps.]

\medskip
\noindent\textbf{19.} Let $A$ be a ring and $\mathfrak p$ a prime ideal of $A$.  Show that every $\mathfrak p$-primary ideal
contains $S_\mathfrak p(0)$, the kernel of the canonical homomorphism $A\longrightarrow A_\mathfrak p$.

Suppose that $A$ satisfies the following condition: for every prime ideal $\mathfrak p$, the
intersection of all $\mathfrak p$-primary ideals of $A$ is equal to $S_\mathfrak p(0)$.  (Noetherian rings
satisfy this condition: see Chapter 10.)  Let $\mathfrak p_1,\ldots,\mathfrak p_n$ be distinct prime ideals,
none of which is a minimal prime ideal of $A$.  Then there exists an ideal $\mathfrak a$ in $A$
whose associated prime ideals are $\mathfrak p_1,\ldots,\mathfrak p_n$.

[Proof by induction on $n$.  The case $n=1$ is trivial (take $\mathfrak a=\mathfrak p_1$).  Suppose
$n>1$ and let $\mathfrak p_n$ be maximal in the set $\{\mathfrak p_1,\ldots,\mathfrak p_n\}$.  By the inductive hypothesis
there exists an ideal $\mathfrak b$ and a minimal primary decomposition $\mathfrak b=\mathfrak q_1\cap\cdots
\cap\mathfrak q_{n-1}$, where each $\mathfrak q_i$ is $\mathfrak p_i$-primary.  If $\mathfrak b\subseteq S_{\mathfrak p_n}(0)$, let $\mathfrak p$ be a minimal prime
ideal of $A$ contained in $\mathfrak p_n$.  Then $S_{\mathfrak p_n}(0)\subseteq S_\mathfrak p(0)$, hence $\mathfrak b\subseteq S_\mathfrak p(0)$.  Taking
radicals and using Exercise 10, we have $\mathfrak p_1\cap\cdots\cap\mathfrak p_{n-1}\subseteq\mathfrak p$, hence some
$\mathfrak p_i\subseteq\mathfrak p$, hence $\mathfrak p_i=\mathfrak p$ since $\mathfrak p$ is minimal.  This is a contradiction since no $\mathfrak p_i$ is
minimal.  Hence $\mathfrak b\nsubseteq S_{\mathfrak p_n}(0)$ and therefore there exists a $\mathfrak p_n$-primary ideal $\mathfrak q_n$
such that $\mathfrak b\nsubseteq\mathfrak q_n$.  Show that $\mathfrak a=\mathfrak q_1\cap\cdots\cap\mathfrak q_n$ has the required properties.]

\medskip
\noindent\emph{Primary decomposition of modules}

Practically the whole of this chapter can be transposed to the context of
modules over a ring $A$.  The following exercises indicate how this is done.

\medskip
\noindent\textbf{20.} Let $M$ be a fixed $A$-module, $N$ a submodule of $M$.  The \emph{radical of $N$ in $M$} is
defined to be
\[
r_M(N)=\{x\in A:x^qM\subseteq N\text{ for some }q>0\}.
\]
Show that $r_M(N)=r(N:M)=r(\operatorname{Ann}(M/N))$.  In particular, $r_M(N)$ is an
ideal.

State and prove the formulas for $r_M$ analogous to \amref{1.13}.

\medskip
\noindent\textbf{21.} An element $x\in A$ defines an endomorphism $\phi_x$ of $M$, namely $m\longmapsto xm$.  The
element $x$ is said to be a \emph{zero-divisor} (resp. \emph{nilpotent}) in $M$ if $\phi_x$ is not injective
(resp. is nilpotent).  A submodule $Q$ of $M$ is \emph{primary in $M$} if $Q\ne M$ and every
zero-divisor in $M/Q$ is nilpotent.

Show that if $Q$ is primary in $M$, then $(Q:M)$ is a primary ideal and hence
$r_M(Q)$ is a prime ideal $\mathfrak p$.  We say that $Q$ is $\mathfrak p$-primary (in $M$).

Prove the analogues of \amref{4.3} and \amref{4.4}.

\medskip
\noindent\textbf{22.} A \emph{primary decomposition of $N$ in $M$} is a representation of $N$ as an intersection
\[
N=Q_1\cap\cdots\cap Q_n
\]
of primary submodules of $M$; it is a \emph{minimal primary decomposition} if the ideals
$\mathfrak p_i=r_M(Q_i)$ are all distinct and if none of the components $Q_i$ can be omitted
from the intersection, that is if $Q_i\nsupseteq\bigcap_{j\ne i}Q_j$ $(1\leq i\leq n)$.

Prove the analogue of \amref{4.5}, that the prime ideals $\mathfrak p_i$ depend only on $N$
(and $M$).  They are called the \emph{prime ideals belonging to $N$ in $M$}.  Show that they
are also the prime ideals belonging to 0 in $M/N$.

\medskip
\noindent\textbf{23.} State and prove the analogues of \amref{4.6}--\amref{4.11} inclusive.  (There is no loss of
generality in taking $N=0$.)

\section{Integral Dependence and Valuations}

In classical algebraic geometry curves were frequently studied by projecting
them onto a line and regarding the curve as a (ramified) covering of the line.
This is quite analogous to the relationship between a number field and the
rational field---or rather between their rings of integers---and the common
algebraic feature is the notion of integral dependence.  In this chapter we prove a
number of results about integral dependence.  In particular we prove the theorems of Cohen-Seidenberg (the ``going-up'' and ``going-down'' theorems)
concerning prime ideals in an integral extension.  In the exercises at the end we
discuss the algebro-geometric situation and in particular the Normalization
Lemma.

We also give a brief treatment of valuations.

\subsection{Integral dependence}

Let $B$ be a ring, $A$ a subring of $B$ (so that $1\in A$).  An element $x$ of $B$ is said to
be \emph{integral over $A$} if $x$ is a root of a monic polynomial with coefficients in $A$, that
is if $x$ satisfies an equation of the form
\begin{equation}
x^n+a_1x^{n-1}+\cdots+a_n=0
\tag{1}
\end{equation}
where the $a_i$ are elements of $A$.  Clearly every element of $A$ is integral over $A$.

\medskip
\noindent\textbf{Example 5.0.\amtarget{5.0}} $A=\mathbb Z$, $B=\mathbb Q$.  If a rational number $x=r/s$ is integral over $\mathbb Z$,
where $r,s$ have no common factor, we have from (1)
\[
r^n+a_1r^{n-1}s+\cdots+a_ns^n=0
\]
the $a_i$ being rational integers.  Hence $s$ divides $r^n$, hence $s=\pm1$, hence $x\in\mathbb Z$.

\medskip
\noindent\textit{\textbf{Proposition 5.1.\amtarget{5.1}} The following are equivalent:}

\noindent\textit{i) $x\in B$ is integral over $A$;}

\noindent\textit{ii) $A[x]$ is a finitely generated $A$-module;}

\noindent\textit{iii) $A[x]$ is contained in a subring $C$ of $B$ such that $C$ is a finitely generated
$A$-module;}

\noindent\textit{iv) There exists a faithful $A[x]$-module $M$ which is finitely generated as an
$A$-module.}

\noindent\emph{Proof.}  i) $\Rightarrow$ ii).  From (1) we have
\[
x^{n+r}=-(a_1x^{n+r-1}+\cdots+a_nx^r)
\]
for all $r\geq0$; hence, by induction, all positive powers of $x$ lie in the $A$-module
generated by $1,x,\ldots,x^{n-1}$.  Hence $A[x]$ is generated (as an $A$-module) by
$1,x,\ldots,x^{n-1}$.

ii) $\Rightarrow$ iii).  Take $C=A[x]$.

iii) $\Rightarrow$ iv).  Take $M=C$, which is a faithful $A[x]$-module (since $yC=0\Rightarrow
y\cdot1=0$).

iv) $\Rightarrow$ i).  This follows from \amref{2.4}: take $\phi$ to be multiplication by $x$, and
$\mathfrak a=A$ (we have $xM\subseteq M$ since $M$ is an $A[x]$-module); since $M$ is faithful, we
have $x^n+a_1x^{n-1}+\cdots+a_n=0$ for suitable $a_i\in A$.\proofbox

\medskip
\noindent\textit{\textbf{Corollary 5.2.\amtarget{5.2}} Let $x_i$ $(1\leq i\leq n)$ be elements of $B$, each integral over $A$.
Then the ring $A[x_1,\ldots,x_n]$ is a finitely-generated $A$-module.}

\noindent\emph{Proof.}  By induction on $n$.  The case $n=1$ is part of \amref{5.1}.  Assume $n>1$, let
$A_r=A[x_1,\ldots,x_r]$; then by the inductive hypothesis $A_{n-1}$ is a finitely generated
$A$-module.  $A_n=A_{n-1}[x_n]$ is a finitely generated $A_{n-1}$-module (by the case
$n=1$, since $x_n$ is integral over $A_{n-1}$).  Hence by \amref{2.16} $A_n$ is finitely generated
as an $A$-module.\proofbox

\medskip
\noindent\textit{\textbf{Corollary 5.3.\amtarget{5.3}} The set $C$ of elements of $B$ which are integral over $A$ is a
subring of $B$ containing $A$.}

\noindent\emph{Proof.}  If $x,y\in C$ then $A[x,y]$ is a finitely generated $A$-module by \amref{5.2}.  Hence
$x\pm y$ and $xy$ are integral over $A$, by iii) of \amref{5.1}.\proofbox

The ring $C$ in \amref{5.3} is called the \emph{integral closure} of $A$ in $B$.  If $C=A$, then $A$
is said to be \emph{integrally closed in $B$}.  If $C=B$, the ring $B$ is said to be \emph{integral
over $A$}.

\emph{Remark.}  Let $f:A\longrightarrow B$ be a ring homomorphism, so that $B$ is an $A$-algebra.
Then $f$ is said to be \emph{integral}, and $B$ is said to be an \emph{integral $A$-algebra}, if $B$ is
integral over its subring $f(A)$.  In this terminology, the above results show that
\[
\text{finite type}+\text{integral}=\text{finite}.
\]

\medskip
\noindent\textit{\textbf{Corollary 5.4.\amtarget{5.4}} If $A\subseteq B\subseteq C$ are rings and if $B$ is integral over $A$, and $C$ is
integral over $B$, then $C$ is integral over $A$ (transitivity of integral dependence).}

\noindent\emph{Proof.}  Let $x\in C$, then we have an equation
\[
x^n+b_1x^{n-1}+\cdots+b_n=0\qquad(b_i\in B).
\]
The ring $B'=A[b_1,\ldots,b_n]$ is a finitely generated $A$-module by \amref{5.2}, and $B'[x]$
is a finitely generated $B'$-module (since $x$ is integral over $B'$).  Hence $B'[x]$ is a
finitely generated $A$-module by \amref{2.16} and therefore $x$ is integral over $A$ by iii)
of \amref{5.1}.\proofbox

\medskip
\noindent\textit{\textbf{Corollary 5.5.\amtarget{5.5}} Let $A\subseteq B$ be rings and let $C$ be the integral closure of $A$
in $B$.  Then $C$ is integrally closed in $B$.}

\noindent\emph{Proof.}  Let $x\in B$ be integral over $C$.  By \amref{5.4} $x$ is integral over $A$, hence
$x\in C$.\proofbox

The next proposition shows that integral dependence is preserved on passing
to quotients and to rings of fractions:

\medskip
\noindent\textit{\textbf{Proposition 5.6.\amtarget{5.6}} Let $A\subseteq B$ be rings, $B$ integral over $A$.}

\noindent\textit{i) If $\mathfrak b$ is an ideal of $B$ and $\mathfrak a=\mathfrak b^c=A\cap\mathfrak b$, then $B/\mathfrak b$ is integral over $A/\mathfrak a$.}

\noindent\textit{ii) If $S$ is a multiplicatively closed subset of $A$, then $S^{-1}B$ is integral over
$S^{-1}A$.}

\noindent\emph{Proof.}  i) If $x\in B$ we have, say, $x^n+a_1x^{n-1}+\cdots+a_n=0$, with $a_i\in A$.
Reduce this equation mod. $\mathfrak b$.

ii) Let $x/s\in S^{-1}B$ ($x\in B$, $s\in S$).  Then the equation above gives
\[
(x/s)^n+(a_1/s)(x/s)^{n-1}+\cdots+a_n/s^n=0
\]
which shows that $x/s$ is integral over $S^{-1}A$.\proofbox

\subsection{The going-up theorem}

\medskip
\noindent\textit{\textbf{Proposition 5.7.\amtarget{5.7}} Let $A\subseteq B$ be integral domains, $B$ integral over $A$.  Then $B$
is a field if and only if $A$ is a field.}

\noindent\emph{Proof.}  Suppose $A$ is a field; let $y\in B$, $y\ne0$.  Let
\[
y^n+a_1y^{n-1}+\cdots+a_n=0\qquad(a_i\in A)
\]
be an equation of integral dependence for $y$ of smallest possible degree.  Since $B$
is an integral domain we have $a_n\ne0$, hence $y^{-1}=-a_n^{-1}(y^{n-1}+a_1y^{n-2}+\cdots
+a_{n-1})\in B$.  Hence $B$ is a field.

Conversely, suppose $B$ is a field; let $x\in A$, $x\ne0$.  Then $x^{-1}\in B$, hence is
integral over $A$, so that we have an equation
\[
x^{-m}+a'_1x^{-m+1}+\cdots+a'_m=0\qquad(a'_i\in A).
\]
It follows that $x^{-1}=-(a'_1+a'_2x+\cdots+a'_mx^{m-1})\in A$, hence $A$ is a field.\proofbox

\medskip
\noindent\textit{\textbf{Corollary 5.8.\amtarget{5.8}} Let $A\subseteq B$ be rings, $B$ integral over $A$; let $\mathfrak q$ be a prime ideal
of $B$ and let $\mathfrak p=\mathfrak q^c=\mathfrak q\cap A$.  Then $\mathfrak q$ is maximal if and only if $\mathfrak p$ is maximal.}

\noindent\emph{Proof.}  By \amref{5.6}, $B/\mathfrak q$ is integral over $A/\mathfrak p$, and both these rings are integral
domains.  Now use \amref{5.7}.\proofbox

\medskip
\noindent\textit{\textbf{Corollary 5.9.\amtarget{5.9}} Let $A\subseteq B$ be rings, $B$ integral over $A$; let $\mathfrak q,\mathfrak q'$ be prime
ideals of $B$ such that $\mathfrak q\subseteq\mathfrak q'$ and $\mathfrak q^c=\mathfrak q'{}^c=\mathfrak p$ say.  Then $\mathfrak q=\mathfrak q'$.}

\noindent\emph{Proof.}  By \amref{5.6}, $B_\mathfrak p$ is integral over $A_\mathfrak p$.  Let $\mathfrak m$ be the extension of $\mathfrak p$ in $A_\mathfrak p$ and let
$\mathfrak n,\mathfrak n'$ be the extensions of $\mathfrak q,\mathfrak q'$ respectively in $B_\mathfrak p$.  Then $\mathfrak m$ is the maximal ideal
of $A_\mathfrak p$; $\mathfrak n\subseteq\mathfrak n'$, and $\mathfrak n^c=\mathfrak n'{}^c=\mathfrak m$.  By \amref{5.8} it follows that $\mathfrak n,\mathfrak n'$ are maximal,
hence $\mathfrak n=\mathfrak n'$, hence by \amref{3.11}(iv) $\mathfrak q=\mathfrak q'$.\proofbox

\medskip
\noindent\textit{\textbf{Theorem 5.10.\amtarget{5.10}} Let $A\subseteq B$ be rings, $B$ integral over $A$, and let $\mathfrak p$ be a prime
ideal of $A$.  Then there exists a prime ideal $\mathfrak q$ of $B$ such that $\mathfrak q\cap A=\mathfrak p$.}

\noindent\emph{Proof.}  By \amref{5.6}, $B_\mathfrak p$ is integral over $A_\mathfrak p$, and the diagram
\[
\begin{array}{ccc}
A&\longrightarrow&B\\
{\scriptstyle\alpha}\big\downarrow&&\big\downarrow{\scriptstyle\beta}\\
A_\mathfrak p&\longrightarrow&B_\mathfrak p
\end{array}
\]
(in which the horizontal arrows are injections) is commutative.  Let $\mathfrak n$ be a maximal ideal of $B_\mathfrak p$; then $\mathfrak m=\mathfrak n\cap A_\mathfrak p$ is maximal by \amref{5.8}, hence is the unique
maximal ideal of the local ring $A_\mathfrak p$.  If $\mathfrak q=\beta^{-1}(\mathfrak n)$, then $\mathfrak q$ is prime and we have
$\mathfrak q\cap A=\alpha^{-1}(\mathfrak m)=\mathfrak p$.\proofbox

\medskip
\noindent\textit{\textbf{Theorem 5.11.\amtarget{5.11}} (``Going-up theorem''). Let $A\subseteq B$ be rings, $B$ integral
over $A$; let $\mathfrak p_1\subseteq\cdots\subseteq\mathfrak p_n$ be a chain of prime ideals of $A$ and $\mathfrak q_1\subseteq\cdots\subseteq\mathfrak q_m$
$(m<n)$ a chain of prime ideals of $B$ such that $\mathfrak q_i\cap A=\mathfrak p_i$ $(1\leq i\leq m)$.
Then the chain $\mathfrak q_1\subseteq\cdots\subseteq\mathfrak q_m$ can be extended to a chain $\mathfrak q_1\subseteq\cdots\subseteq\mathfrak q_n$ such
that $\mathfrak q_i\cap A=\mathfrak p_i$ for $1\leq i\leq n$.}

\noindent\emph{Proof.}  By induction we reduce immediately to the case $m=1$, $n=2$.  Let
$\bar A=A/\mathfrak p_1$, $\bar B=B/\mathfrak q_1$; then $\bar A\subseteq\bar B$, and $\bar B$ is integral over $\bar A$ by \amref{5.6}.  Hence, by
\amref{5.10}, there exists a prime ideal $\bar{\mathfrak q}_2$ of $\bar B$ such that $\bar{\mathfrak q}_2\cap\bar A=\bar{\mathfrak p}_2$, the image of $\mathfrak p_2$
in $\bar A$.  Lift back $\bar{\mathfrak q}_2$ to $B$ and we have a prime ideal $\mathfrak q_2$ with the required properties.\proofbox

\subsection{Integrally closed integral domains. The going-down theorem}

Proposition \amref{5.6}(ii) can be sharpened:

\medskip
\noindent\textit{\textbf{Proposition 5.12.\amtarget{5.12}} Let $A\subseteq B$ be rings, $C$ the integral closure of $A$ in $B$.  Let $S$
be a multiplicatively closed subset of $A$.  Then $S^{-1}C$ is the integral closure of
$S^{-1}A$ in $S^{-1}B$.}

\noindent\emph{Proof.}  By \amref{5.6}, $S^{-1}C$ is integral over $S^{-1}A$.  Conversely, if $b/s\in S^{-1}B$ is
integral over $S^{-1}A$, then we have an equation of the form
\[
(b/s)^n+(a_1/s_1)(b/s)^{n-1}+\cdots+a_n/s_n=0
\]
where $a_i\in A$, $s_i\in S$ $(1\leq i\leq n)$.  Let $t=s_1\cdots s_n$ and multiply this equation by
$(st)^n$ throughout.  Then it becomes an equation of integral dependence for $bt$
over $A$.  Hence $bt\in C$ and therefore $b/s=bt/st\in S^{-1}C$.\proofbox

An integral domain is said to be \emph{integrally closed} (without qualification)
if it is integrally closed in its field of fractions.  For example, $\mathbb Z$ is integrally
closed (see \amref{5.0}).  The same argument shows that any unique factorization
domain is integrally closed.  In particular, a polynomial ring $k[x_1,\ldots,x_n]$ over
a field is integrally closed.

Integral closure is a local property:

\medskip
\noindent\textit{\textbf{Proposition 5.13.\amtarget{5.13}} Let $A$ be an integral domain.  Then the following are
equivalent:}

\noindent\textit{i) $A$ is integrally closed;}

\noindent\textit{ii) $A_\mathfrak p$ is integrally closed, for each prime ideal $\mathfrak p$;}

\noindent\textit{iii) $A_\mathfrak m$ is integrally closed, for each maximal ideal $\mathfrak m$.}

\noindent\emph{Proof.}  Let $K$ be the field of fractions of $A$, let $C$ be the integral closure of $A$ in $K$,
and let $f:A\longrightarrow C$ be the identity mapping of $A$ into $C$.  Then $A$ is integrally
closed $\Longleftrightarrow f$ is surjective, and by \amref{5.12} $A_\mathfrak p$ (resp. $A_\mathfrak m$) is integrally closed $\Longleftrightarrow f_\mathfrak p$
(resp. $f_\mathfrak m$) is surjective.  Now use \amref{3.9}.\proofbox

Let $A\subseteq B$ be rings and let $\mathfrak a$ be an ideal of $A$.  An element of $B$ is said to be
\emph{integral over $\mathfrak a$} if it satisfies an equation of integral dependence over $A$ in which
all the coefficients lie in $\mathfrak a$.  The \emph{integral closure of $\mathfrak a$ in $B$} is the set of all elements
of $B$ which are integral over $\mathfrak a$.

\medskip
\noindent\textit{\textbf{Lemma 5.14.\amtarget{5.14}} Let $C$ be the integral closure of $A$ in $B$ and let $\mathfrak a^e$ denote the
extension of $\mathfrak a$ in $C$.  Then the integral closure of $\mathfrak a$ in $B$ is the radical of $\mathfrak a^e$
(and is therefore closed under addition and multiplication).}

\noindent\emph{Proof.}  If $x\in B$ is integral over $\mathfrak a$, we have an equation of the form
\[
x^n+a_1x^{n-1}+\cdots+a_n=0
\]
with $a_1,\ldots,a_n$ in $\mathfrak a$.  Hence $x\in C$ and $x^n\in\mathfrak a^e$, that is $x\in r(\mathfrak a^e)$.  Conversely, if
$x\in r(\mathfrak a^e)$ then $x^n=\sum a_ix_i$ for some $n>0$, where the $a_i$ are elements of $\mathfrak a$ and the
$x_i$ are elements of $C$.  Since each $x_i$ is integral over $A$ it follows from \amref{5.2} that
$M=A[x_1,\ldots,x_n]$ is a finitely generated $A$-module, and we have $x^nM\subseteq\mathfrak aM$.
Hence by \amref{2.4} (taking $\phi$ there to be multiplication by $x^n$) we see that $x^n$ is
integral over $\mathfrak a$, hence $x$ is integral over $\mathfrak a$.\proofbox

\medskip
\noindent\textit{\textbf{Proposition 5.15.\amtarget{5.15}} Let $A\subseteq B$ be integral domains, $A$ integrally closed, and
let $x\in B$ be integral over an ideal $\mathfrak a$ of $A$.  Then $x$ is algebraic over the field of
fractions $K$ of $A$, and if its minimal polynomial over $K$ is $t^n+a_1t^{n-1}+\cdots
+a_n$, then $a_1,\ldots,a_n$ lie in $r(\mathfrak a)$.}

\noindent\emph{Proof.}  Clearly $x$ is algebraic over $K$.  Let $L$ be an extension field of $K$ which
contains all the conjugates $x_1,\ldots,x_n$ of $x$.  Each $x_i$ satisfies the same equation
of integral dependence as $x$ does, hence each $x_i$ is integral over $\mathfrak a$.  The coefficients of the minimal polynomial of $x$ over $K$ are polynomials in the $x_i$, hence
by \amref{5.14} are integral over $\mathfrak a$.  Since $A$ is integrally closed, they must lie in $r(\mathfrak a)$, by
\amref{5.14} again.\proofbox
